% Generated from validated Article Draft Result. Do not hand-edit; revise the structured draft instead.
\section{AgentはUIを出て実行系になる — Harness・Context・Tool・Browser/Computer Use}
\label{pkg:agent}
\noindent\textbf{2025年後半、Agent化は会話UIの拡張ではなく、browser、coding surface、tool/API、評価基盤を含む実行系の設計問題へ広がった。}\autocite{src-7a7213b68631,src-b074ccdb2d9e}

7月17日のChatGPT agentはresearch、browser automation、code/tool executionをuser-facing productへまとめ、同日system cardも公開した。action capabilityとcontrol boundaryが同じrelease packageで前景化した。\autocite{src-7a7213b68631,src-b074ccdb2d9e}


AnthropicはSonnet 4.5、Haiku 4.5、Opus 4.5をAPIやcoding/agent platformと結び付けて展開した。newsroomはliving indexなので、2025年後半のchronologyはitem-level公式ページで再構成する必要がある。\autocite{src-f6aa679babd7}


Gemini APIはvideo/image、robotics、Computer Use、Gemini 3、Interactions API、Deep Researchへsurfaceを拡大した。preview、GA、deprecation、shutdownは独立したservice lifecycle eventであり、upstream model announcementへ畳み込まない。\autocite{src-a24fd97eb3d0}


10月のAgentKit、Evals、RFT for agentsは、agentの構築だけでなく評価と改善を反復するengineering layerを示した。\autocite{src-0661418088b8}


9月のCodex upgradesはterminal、IDE、web、mobileへcoding-agent surfaceを広げた。modelだけでなくcontext handoff、tool execution、human reviewを含むexecution stackが実用能力を左右する。\autocite{src-6deeebce1719}
